\documentclass[11pt,a4paper]{article}
\usepackage[T1]{fontenc}
\usepackage{lmodern,amsmath,amssymb,amsthm,mathtools,microtype}
\usepackage[margin=27mm]{geometry}
\usepackage[colorlinks=true,linkcolor=blue,urlcolor=blue]{hyperref}
\newtheorem{proposition}{Proposition}
\newcommand{\Q}{\mathbb Q}
\title{Failure of a two-square mixed-power height bound}
\author{}\date{17 November 2026}
\begin{document}\maketitle

For a reduced rational number $a/b$ with $b>0$, write
$h(a/b)=\log\max(|a|,b)$.

\begin{proposition}\label{prop:failure}
For every $d=2^r$, $r\ge1$, there are no constants $C,K$ such
that all rational solutions of
\[
                         t=M^d+z_1^2+z_2^2>0
\]
satisfy $h(M)\le C h(t)+K$. The failure persists after requiring
$M$ to exceed an arbitrary fixed positive real lower bound.
For $d\ge4$, both square variables may be required to be nonzero.
\end{proposition}
\begin{proof}
Let $N\ge2$ be an integer and put
\[
 a_N=\frac{N^2-1}{N^2+1},\qquad
 b_N=\frac{2N}{N^2+1}.
\]
Then $a_N^2+b_N^2=1$. The identity
\[
 1-a_N^{2^r}
  =(1-a_N^2)\prod_{j=1}^{r-1}(1+a_N^{2^j})
\]
is a Gaussian norm: take the norm of
\[
 b_N\prod_{j=1}^{r-1}(1+i a_N^{2^{j-1}}).
\]
Its real and imaginary parts are rational; call them $u_N,v_N$.
Thus $1-a_N^d=u_N^2+v_N^2$.

Fix an integer $L>0$ and set
\[
 M_N=L a_N,\qquad t=L^d,\qquad
 z_{1,N}=L^{d/2}u_N,\qquad z_{2,N}=L^{d/2}v_N.
\]
These satisfy the displayed equation with $t$ fixed. The reduced
denominator of $M_N$ is at least $(N^2+1)/(2L)$, since
$\gcd(N^2-1,N^2+1)\le2$. Hence $h(M_N)\to\infty$.
Also $3L/5\le M_N<L$. Taking $L$ sufficiently large gives any
prescribed positive lower bound. For $d\ge4$, the two Gaussian
components are nonzero rational functions of $a_N$, so excluding
finitely many values of $N$ makes both nonzero.
\end{proof}

For the admissible exponent $d=8$, the formula is particularly short:
\begin{equation}\label{eq:eight}
 \begin{gathered}
 a=\frac{N^2-1}{N^2+1},\qquad b=\frac{2N}{N^2+1},\qquad
 M=La,\qquad t=L^8,\\
 z_1=L^4 b(1-a^3),\qquad z_2=L^4 b(a+a^2).
 \end{gathered}
\end{equation}
Both $z_1,z_2$ are positive for $N\ge2$, and direct expansion gives
\[
 (1-a^3)^2+(a+a^2)^2=(1+a^2)(1+a^4),
 \qquad M^8+z_1^2+z_2^2=L^8.
\]
The exact identities and sample rational specializations are verified
by the accompanying checker,
\nolinkurl{check_two_square_height.py}.

\section*{Consequences for a single-detector proposal}

The one-square proof of
\[
              h(M)\le C_{\kappa,d}h(M^d+z^2)+O(1)
\]
uses a three-term coprime identity $A^d+B^2=c'D^d$ after clearing
denominators. Adding a second independent square gives a fourth term;
the same application of $abc$ is not available. Proposition
\ref{prop:failure} shows that no argument, conditional on $abc$ or
otherwise, can supply its proposed unrestricted two-square analogue
for the powers of two, including $d=8$ and $d=16$.

The counterexample also shows why the real inequality $t\ge M^d$
does not control logarithmic rational height. The values of $M$ can
stay in a fixed large real interval while their reduced denominators
grow without bound.

In the intended detector source there is an additional restriction
\[
               M=(X^2+B)^e+T(n)^e.
\]
The family in this note does not assert that restriction, and does
not supply a counterexample to a fully specified detector source.
A height theorem for this special inner form would require a new
proof. In particular, rational sum-of-two-squares preparation at
selected primes does not establish soundness of a twelve-coordinate
source by the existing one-square argument. Exponents with an odd
factor, such as $d=6$, are not ruled out by this particular family;
the original proof still does not transfer to them.
\end{document}
